\section{The Six Frozen Objects}
\label{sec:candidates}

\begin{table}[t]
\centering
\caption{Frozen same-object gate outcomes. S denotes local structural support, P denotes weak/partial/formal evidence, and F denotes failure. Route eligibility is the sequential conjunction on one row; all six Route-B flags are locked.}
\label{tab:candidate-matrix}
\small
\setlength{\tabcolsep}{4.2pt}
\resizebox{\textwidth}{!}{%
\begin{tabular}{lcccccccl}
\toprule
Candidate & A0 & A1 & A2 & A3 & A4 & First block & Evidence & Status \\
\midrule
SD-C01 & P & S & S & F & F & A0 & Proved & Rejected \\
SD-C02 & F & F & S & F & F & A0 & Modeling choice & Rejected \\
SD-C03 & F & P & F & F & F & A0 & Modeling choice & Rejected \\
SD-C04 & P & S & S & P & P & A0 & Proved & Exploratory \\
SD-C05 & S & F & F & F & F & A1 & Proved & Exploratory \\
SD-C06 & S & F & F & P & F & A1 & Not testable & Exploratory \\
\bottomrule
\end{tabular}%
}
\end{table}


\Cref{tab:candidate-matrix} is generated from the evaluator records rather
than transcribed into the manuscript.  The label \emph{supported} means
supported for the candidate's own declared species and convention; it does
not mean agreement with the completed Riemann divisor.

\subsection{\candidate{01}: finite-state arithmetic skeleton}

For the full \(q\)-shift, primitive necklaces of length \(n\) have count
\[
  N_q(n)=\frac{1}{n}\sum_{d\mid n}\mu(d)q^{n/d},
\]
which is also the count of monic irreducible degree-\(n\) polynomials over
\(\F_q\).  The repetitions are exact, and after the normalization
\(u=q^{-s}\) the inverse zeta is
\[
  D_q(s)=1-q^{1-s}.
\]
This is a complete function-field arithmetic skeleton, so A1 and A2 are
supported for that species.  It does not provide a canonical rational prime
\(p\mapsto\gamma_p\), and \cref{thm:finite-memory} blocks the broader
finite-memory comparison class at A3.  The candidate is rejected rather than
enlarged by adding more finite states.

\subsection{\candidate{02}: squarefree admissibility}

The phase space is the two-sided binary shift whose support misses a residue
class modulo \(p^2\) for every rational prime \(p\).  This is genuine
arithmetic structure, but the primes occur explicitly in the grammar, so the
strict emergence reading of A0 fails.  More decisively,
\cref{prop:squarefree} leaves only \(0^\Z\) as a periodic point and gives
\(D_{\mathrm{AM}}(z)=1-z\).  The candidate illustrates why rich aperiodic
language structure and an arithmetic definition do not imply a useful
primitive-orbit ledger.

\subsection{\candidate{03}: shared-base renewal inverse design}

The renewal graph has a distinguished base vertex \(b\).  For each \(n\ge1\),
it has a first-return atom \(r_n\): a length-\(n\) directed path from \(b\)
back to \(b\), with internal vertices disjoint from all other atoms and
aggregate complex weight \(a_nz^n\).  At \(b\), any atom may follow any other atom, and
the first-return factorization is unique.  Hence
\[
  F(z)=\sum_{n\ge1}a_nz^n,\qquad
  Z_{\mathrm{ren}}(z)=\frac{1}{1-F(z)},\qquad
  D_{\mathrm{ren}}(z)=1-F(z).
\]
Primitive cycles are cyclic necklaces in the return atoms.  This gives a weak
primitive expansion, but three exact failures intervene.  The
weights contain all analytic information
(\cref{prop:inverse-design}); atom concatenation creates unwanted mixed
primitive words; and finite unitary phases cannot remove their factors
(\cref{prop:mixed}).  The same degree-12 procedure reconstructs on-circle,
off-circle, and generic controls exactly.  The outcome is therefore
\evidence{stop scoped / proves too much}; no Riemann target was fitted after
that stop.

\subsection{\candidate{04}: Gauss words and the Mayer operator}

For a continued-fraction digit \(n\), write
\[
  \phi_n(z)=\frac{1}{n+z},\qquad
  A_n=\begin{pmatrix}0&1\\1&n\end{pmatrix}.
\]
For \(w=(n_1,\ldots,n_m)\), take the ordered composition
\(\phi_w=\phi_{n_1}\circ\cdots\circ\phi_{n_m}\) and its matching Möbius
product \(M_w=A_{n_1}\cdots A_{n_m}\).  (The experiment uses a fixed
conjugate matrix convention, which preserves traces, eigenvalues, and roofs.)
A primitive digit necklace determines an attracting fixed point \(x_w\), and
the derivative roof satisfies
\[
  T_w=-\log\lvert\phi_w'(x_w)\rvert
      =2\log\lambda_+(M_w),\qquad T_{v^r}=rT_v.
\]
Let
\[
  D=\{z\in\C:\lvert z-1\rvert<3/2\},
\]
and let \(A_\infty(D)\) be the Banach space of functions holomorphic on
\(D\), continuous on its closure, with the supremum norm.  For
\(\RePart(s)>1/2\), Mayer realizes
\[
  (\mathcal L_s f)(z)=\sum_{n\ge1}(z+n)^{-2s}
  f\!\left(\frac1{z+n}\right)
\]
as a nuclear operator of order zero on \(A_\infty(D)\) (and on a
corresponding Hardy-space realization).  In that source domain its Fredholm
determinants satisfy
\[
  D_{\mathrm{MG}}(s)
  =\det(I-\mathcal L_s^2)
  =\det(I-\mathcal L_s)\det(I+\mathcal L_s)
\]
\citep{mayer1990gauss,mayer1991selberg}.  These are genuine A1--A2 successes
and lie outside \cref{thm:finite-memory}.

The mismatch is not analytic naturalness but arithmetic species.  Primitive
continued fractions encode quadratic irrationals or hyperbolic modular
classes, not rational primes.  At the largest finite cutoff, many distinct
non-reversal classes share a matrix trace, so trace is not an injective
arithmetic label.  That collision does not rule out every future map, but no
map frozen here passes A0.  A modular-surface interpretation is not imported
to repair the same-object symbolic audit.

\subsection{\candidate{05}: wheel-sieve level shift}

The recursive least-coprime rule supplies the strongest endogenous
rational-prime origin among the six objects.  With \(Q_0=1\) and
\(q_1=Q_1=2\), level \(k\) has vertices
given by the unit residues
\[
  R_k=\{0\le r<Q_k:\gcd(r,Q_k)=1\}.
\]
The vertex \((k,r)\) has an edge to each unit lift
\((k+1,r+jQ_k)\), with \(0\le j<q_{k+1}\); exactly one lift is deleted
because it is divisible by \(q_{k+1}\).  If \(X_k\) is the set of
one-sided path tails beginning at level \(k\), then deleting the first edge
defines \(\sigma:\bigsqcup_{k\ge0}X_k\to\bigsqcup_{k\ge0}X_k\) with
\(\sigma(X_k)\subset X_{k+1}\).

\Cref{prop:wheel} proves both \(q_k=p_k\) and the intrinsic scale increment
\(\log p_{k+1}\), so A0 is supported without a prime table.  The same
definition shows why this is not yet a zeta candidate: every edge raises the
level, no positive-period orbit exists, and the formal determinant is \(1\).
The arithmetic coordinate cannot be transferred to \candidate{04}'s
determinant.  A stationary factor or recoding would be a new source lock, not
an A1 result for the frozen DAG.

\subsection{\candidate{06}: Knauf's arithmetic recursion}

On finite binary words set \(h(\varnothing)=1\) and, for a prefix
\(u\in\{0,1\}^{k-1}\), define
\[
  h(u0)=h(u),\qquad h(u1)=h(u)+h(\bar u),
\]
where \(\bar u\) is the bitwise complement.  Let \(h_k\) be the restriction
to \(\{0,1\}^k\), and group configurations by
\[
  \varphi_k(n)=\#\{\sigma\in\{0,1\}^k:h_k(\sigma)=n\}.
\]
Then
\[
  Z_k(s)=\sum_n\varphi_k(n)n^{-s}.
\]
The source construction defines \(h\) on the finite-support direct union,
identifies its full multiplicity with Euler's totient, and proves
\(0\le\varphi_k(n)\le\varphi(n)\) together with absolute convergence
\[
  \lim_{k\to\infty}Z_k(s)
  =\sum_{n\ge1}\frac{\varphi(n)}{n^s}
  =\frac{\zeta(s-1)}{\zeta(s)},\qquad \RePart(s)>2
\]
\citep{knauf2013adelic,knauf1998spin,knauf1999erratum}.  No universal
finite-\(k\) equality range is assumed here: the separately reported
\(k=22\) audit certifies only the finite statement
\(\varphi_{22}(n)=\varphi(n)\) for \(1\le n\le23\).  The exact limiting
quotient is the strongest analytic-arithmetic collision in the audit.

The result does not supply A1 or A2.  The finite configurations and their
Dirichlet partition sum have no frozen orientation, repetition, monodromy, or
complete primitive-cycle map.  The Liouville-weighted refinement uses an
extra arithmetic sign unless a symbolic symmetry derives it.  Finite-depth
values at \(\RePart(s)\le2\) remain continuation or boundary benchmarks, not
evidence for the open signed-convergence region.  In particular, the
candidate supplies neither a completed-\(\xi\) Fredholm divisor nor an
operator lift.
